% ---
% title: Rating bonds
% tags: finance, multinomial
% ---

\begin{problem}
	\begin{context}
		A \emph{credit rating} (e.g.\ AAA, A, BBB, \ldots, given by
		agencies such as S\&P, Moody's, Fitch) summarizes how likely a borrower is
		to repay. Here an analyst uses three classes A, B, C (best to worst).
	\end{context}
	A credit analyst must rate 10 bonds: exactly 4 get rating A, 3 get
	rating B and 3 get rating C. How many rating assignments are possible?
\end{problem}

\begin{slides}
	\small
	\exstatement
	\pause
	\sol Three labeled groups of sizes 4, 3, 3:
	\[ \binom{10}{4,3,3} = \frac{10!}{4!\,3!\,3!}
		= \binom{10}{4}\binom63 = 210\cdot 20 = 4200. \]
\end{slides}

\begin{solution}
	\idea The 10 bonds are different; we divide them into three labeled groups
	(A, B, C) of sizes 4, 3, 3.
	\[ \binom{10}{4,3,3}=\frac{10!}{4!\,3!\,3!}=\binom{10}{4}\binom63\binom33
		=210\cdot20\cdot1=4200. \]
	\begin{answer}
		$4200$
	\end{answer}
	\begin{remark}
		Groups B and C have the same size, but they are \emph{not} interchangeable
		(rating B is not rating C), so we do not divide by $2!$.
	\end{remark}
\end{solution}
